% id: 1030
% book: RPM 공통수학2 · 09 유리함수
% section: 시험에 꼭 나오는 문제
% figure: none
% answer: ④
% answer_source: 답지
% variant_level: 0 (원본 전사)
% 이 파일은 build.mjs 가 items.json 에서 만든다. 고칠 때는 items.json 을 고친다.
\hspace*{2mm}\begin{minipage}[t]{\dimexpr\linewidth-2mm\relax}\vspace{0pt}\smash{\raisebox{1.5mm}{\dmkichul{RPM 공통수학2 1030번 · 152쪽}}}
다음 중 함수 $y=\dfrac{k}{x}\ (k\ne 0)$에 대한 설명으로 옳지 \underline{않은} 것은?
\par\medskip\par\noindent\hangindent=1.4em\hangafter=1 {\small ①}\ 그래프의 점근선의 방정식은 $x=0$, $y=0$이다.\par\noindent\hangindent=1.4em\hangafter=1 {\small ②}\ 그래프는 직선 $y=-x$에 대하여 대칭이다.\par\noindent\hangindent=1.4em\hangafter=1 {\small ③}\ 정의역과 치역이 같다.\par\noindent\hangindent=1.4em\hangafter=1 {\small ④}\ $|k|$의 값이 클수록 그래프는 원점에 가까워진다.\par\noindent\hangindent=1.4em\hangafter=1 {\small ⑤}\ $k<0$이면 그래프는 제$2$사분면과 제$4$사분면을 지난다.\par\medskip\end{minipage}
